\cotefile{157-5}
\scanpage{133}
% printed form
OU ENCORE APPELEZ 'SVP RESEAU' (67) 54 41 \unsure{78|75|73}.

\unsure{ET DE L'ORDINATEUR|ET DE L'ORDINATEUR :}

LUNDI 3 NOVEMBRE DE \unsure{7|2} A \unsure{8|9} H

\unsure{MARDI 4 NOVEMBRE DE 7 A 9 H}

\handwritten{\deleted{1 Rappels sur les couples de Quillen \unsure{p.|P.} 1-12}}

\handwritten{\deleted{2 Rappels Triples de Quillen, triples de Q. spéciaux p. 13-\unsure{32|33} (th. 2.3 p. \unsure{19|14}, th. 2.6. p. 25, th. 2.9. p. \unsure{28|29})}}

\handwritten{\deleted{3. Catégories à \unsure{cofibrations|fibrations} \quad p. 33-50}}

\handwritten{4. Paratopos et \unsure{l'opération|Opération} $M\unsure{\boxtimes|\boxed{\rtimes}}N$ \quad p. 50-}

\handwritten{4.1. Pseudo topos \quad p. 50-57}

\handwritten{4.2. Préliminaires sur $M\boxtimes N$ \quad p. 57-\unsure{69|68|64}}

\handwritten{\deleted{(th. 4.3.6 p. 70-71)}}

\handwritten{4.3. \deleted{$\operatorname{Ind}_{\pi} C$ : \unsure{stabilité} par $\underset{\longrightarrow}{\lim}$} \inserted{$M\boxtimes N$ et \unsure{exigences|esquisse} \unsure{sur|de} l'hyp \unsure{des|du} ps. topos} \quad p. \unsure{69|68|64}-\unsure{78|79} (th. 4.3.6 p. 70-71)}

\handwritten{4.4. $\operatorname{Ind}_{\pi} C$ (stabilité par $\underset{\longrightarrow}{\lim}$. \quad p. \unsure{79|78}-106 (th. 4.4.13 p. 95-96)}

% printed form
GUTS OUTPUT FOR PREDONT DELIVER AT 0000 - PREDONT.MEFJCLAP

GUTS OUTPUT FOR PREDONT DELIVER AT 0000 - PREDONT.MEFJCLAP

GUTS OUTPUT FOR PREDONT DELIVER AT 0000 - PREDONT.MEFJCLAP

GUTS OUTPUT FOR PREDONT DELIVER AT 0000 - PREDONT.MEFJCLAP

GUTS OUTPUT FOR PREDONT DELIVER AT 0000 - PREDONT.MEFJCLAP

GUTS OUTPUT FOR PREDONT DELIVER AT 0000 - PREDONT.MEFJCLAP

GUTS OUTPUT FOR PREDONT DELIVER AT 0000 - PREDONT.MEFJCLAP

GUTS OUTPUT FOR PREDONT DELIVER AT 0000 - PREDONT.MEFJCLAP

GUTS OUTPUT FOR PREDONT DELIVER AT 0000 - PREDONT.MEFJCLAP

GUTS OUTPUT FOR PREDONT DELIVER AT 0000 - PREDONT.MEFJCLAP

GUTS OUTPUT FOR PREDONT DELIVER AT 0000 - PREDONT.MEFJCLAP

GUTS OUTPUT FOR PREDONT DELIVER AT 0000 - PREDONT.MEFJCLAP

GUTS OUTPUT FOR PREDONT DELIVER AT 0000 - PREDONT.MEFJCLAP

GUTS OUTPUT FOR PREDONT DELIVER AT 0000 - PREDONT.MEFJCLAP

GUTS OUTPUT FOR PREDONT DELIVER AT 0000 - PREDONT.MEFJCLAP

GUTS OUTPUT FOR PREDONT DELIVER AT 0000 - PREDONT.MEFJCLAP

GUTS OUTPUT FOR PREDONT DELIVER AT 0000 - PREDONT.MEFJCLAP

GUTS OUTPUT FOR PREDONT DELIVER AT 0000 - PREDONT.MEFJCLAP

GUTS OUTPUT FOR PREDONT DELIVER AT 0000 - PREDONT.MEFJCLAP

GUTS OUTPUT FOR PREDONT DELIVER AT 0000 - PREDONT.MEFJCLAP

GUTS OUTPUT FOR PREDONT DELIVER AT 0000 - PREDONT.MEFJCLAP

GUTS OUTPUT FOR PREDONT DELIVER AT 0000 - PREDONT.MEFJCLAP

GUTS OUTPUT FOR PREDONT DELIVER AT 0000 - PREDONT.MEFJCLAP

GUTS OUTPUT FOR PREDONT DELIVER AT 0000 - PREDONT.MEFJCLAP

GUTS OUTPUT FOR PREDONT DELIVER AT 0000 - PREDONT.MEFJCLAP

GUTS OUTPUT FOR PREDONT DELIVER AT 0000 - PREDONT.MEFJCLAP

GUTS OUTPUT FOR PREDONT DELIVER AT 0000 - PREDONT.MEFJCLAP

GUTS OUTPUT FOR PREDONT DELIVER AT 0000 - PREDONT.MEFJCLAP

GUTS OUTPUT FOR PREDONT DELIVER AT 0000 - PREDONT.MEFJCLAP

GUTS OUTPUT FOR PREDONT DELIVER AT 0000 - PREDONT.MEFJCLAP

GUTS OUTPUT FOR PREDONT DELIVER AT 0000 - PREDONT.MEFJCLAP

GUTS OUTPUT FOR PREDONT DELIVER AT 0000 - PREDONT.MEFJCLAP

\handwritten{4.5. $\operatorname{Ind}_{\pi}C$ (\unsure{filtration cardinale}) \quad p. 106-139 (Cf \unsure{4.5.6.|4.5.6.3.6.} p. 116)}

\handwritten{4.6. \unsure{Dém.|Dim.} de (\uline{Caractérisation des paratopos}) \quad p. 139-\deleted{\unsure{165|167}}160 th. récapitulatif \fbox{\fbox{4.6.3}} p. \unsure{148|143}-150)}

\handwritten{4.7. \deleted{Stabilité des paratopos} \quad \deleted{p. 160-1} $\operatorname{Ind}_{\pi}(C)$ ($C$ quelconque) et filtrations \unsure{cardinales|ordinales} \unsure{génératrices|généralisées} \quad p. \unsure{161|162}-187,}

\handwritten{4.8. Catégories \inserted{$\unsure{U|\mathcal U}$-}accessibles. \quad p. 187-216}

\handwritten{4.9. Le théorème des représentations indicielles \quad p \unsure{217|227}-242 (th 4.9.3 p. 219)}

\handwritten{4.10. Dém. et généralisations du th. des représentations indicielles \quad p. 243-267 (th. 4.10.\unsure{1|4} p. 244)}

\handwritten{4.11. Application aux limites des $\operatorname{Ind}_{\pi}C$ \quad p. 267-273}

\handwritten{4.12. Foncteurs accessibles \quad p. 273-280}

\handwritten{4.13. Parties accessibles de $\operatorname{Ob}M$ \quad p. \unsure{281|289}-369 (\unsure{cf :} 4.13.\unsure{32|22} p. \unsure{342|372},) (4.13.31 p. 350)}

\handwritten{4.14. Critères d'accessibilité pour \deleted{\unsure{les|la}} une catégorie \quad p. 370-424}

\handwritten{4.15. Catégories accessibles sur une autre \quad p. 425-463}

\handwritten{4.16. Exemples \unsure{et|de} questions \quad \deleted{p. 464-\unsure{486|485} \deleted{\unsure{485|488}}} p. 464-491}

\handwritten{4.17. Retour sur $M\boxtimes N$\ldots{}- \quad p. 492-}
% archive sheet number: 132
\scanpage{134}
\mspage{50}

\deleted{\unsure{ça|ça.} ce sont les exemples qui manquent. La clef d'une compréhension semble bien être dans les questions de \unsure{fonctorialité|fonctorisation} entre ensembles $\Phi\longleftrightarrow\Psi$. Il faudrait en faire une étude systématique, mais je recule devant la perspective d'une digression qui me prendrait peut-être encore des semaines !}

4. \inserted{\uline{Paratopos et}} \uline{L'opération} $M\boxtimes N$.

4.1. \inserted{\uline{Pseudo-topos.} \uline{Pb. du \unsure{produit|premier}}}

J'appelle \inserted{\uline{Catégorie}} \uline{pseudo-$\unsure{U|\mathcal U}$-topos}, ou simplement \uline{pseudo-topos}, une \inserted{$U$-}catégorie $M$ stable par $U$-limites \deleted{projectives}\inserted{inductives}, et qui admet une \uline{petite} sous-catégorie $M_0$ \deleted{stable par} \uline{strictement génératrice} (i.e. génératrice par \unsure{épimorphismes} stricts). Alors $M$ admet aussi \unsure{ses|des} $U$-limites projectives \inserted{(SGA 4 I 8.\unsure{12|11}.\unsure{7|3})}, \deleted{[Mais} l'exemple de (\unsure{Gr|An|Cn}), cf loc cit 8.12.9, montre que $M^{\mathrm{o}}$ n'est pas forcément, pour autant, un pseudo-$U$-topos, faute à $M$ d'admettre une petite \deleted{\unsure{famille}} sous-catégorie strictement \uline{co}génératrice. Mais le th. 2.6. donne des conditions sur $M$ (\unsure{surtout} du type \unsure{exactitude})
% archive sheet number: 133
\scanpage{135}
\mspage{51}

pour que $M$ admette assez d'injectifs -- ce qui implique alors l'existence d'une petite sous-catégorie strictement \uline{co}génératrice, donc que $M^{\mathrm{o}}$ est également un pseudo-topos.$^{*}$

% slanted
\margin{C'est le cas notamment si $M$ est un \uline{topos}.}

Je \unsure{visualiserai|visualise} les \inserted{catégories} pseudo-topos comme étant la catégorie des ``faisceaux'' sur un objet $\unsure{X|\mathcal X}$, visualisé comme une sorte \unsure{$^{*}$|``} d'``espace'', que je distingue de la catégorie $M=\operatorname{Faisc}(X)$ des faisceaux sur $X$. \deleted{Les} C'est $X$ qui sera pour moi le ``pseudo-topos''. Les pseudo-topos forment une 2-catégorie, la variance des $X$ étant opposée à celle des \deleted{\unsure{faisc}} $\operatorname{Fais}(X)=M$. \deleted{$\underline{\operatorname{Hom}}(X,Y)={}$} De façon précise, si $X$ correspond à $M$, $\unsure{Y|\mathcal Y}$ à $N$, un morphisme $/f/$ de $X$ dans $Y$ sera donné par un foncteur

% slanted
\margin{\unsure{Dire|D'où} ici que les \deleted{\unsure{faisceaux sur}} objets de $M$ sont visualisés comme des \deleted{\unsure{obj.}} ``\unsure{sections|objets}'' généralisés de $X$. On a $M\xrightarrow{\approx}\underline{\operatorname{Hom}}^{!}(M^{\mathrm{o}},\mathrm{Ens})$, ici $M$ s'\unsure{interprète|réinterprète} bel et bien comme les ``faisceaux'' sur cette catégorie $M$ d'\unsure{ensembl.|arguments}\ldots{} cf.}

\[
f^{*}:\deleted{\unsure{M}}N\longrightarrow M
\tag*{(4.1.1)}
\]

en sens inverse, commutant \deleted{\unsure{à}} aux \uline{lim} inductives quelconques. [\uline{NB} si $X,Y$ sont $=$ des topos, i.e. $M,N$ des catégories-topos, \deleted{on exige} et si on veut que $f$ soit un morphisme de topos, on \uline{exige} de plus que $f^{*}$ soit de plus \uline{exact à g.} On dira qu'un \uline{morphisme de pseudo-topos est} \uuline{exact} si de plus $f^{*}$ est exact à g.]
% archive sheet number: 134
\scanpage{136}
\mspage{52}

En vertu du dual de loc cit 8.12.7, cela \deleted{implique}\inserted{\unsure{signifie en fait|signifie au fond}} que $f^{*}$ admette un adjoint à droite $f_{*}$

\[
f_{*}:M\unsure{\longrightarrow|\text{---}}N,
\tag*{(4.1.2)}
\]

lequel commutera donc aux petites $\underset{\longleftarrow}{\lim}$. La donnée de $f_{*}$ satisfaisant à cette propriété \unsure{détermine|déterminera} \unsure{donc|son} $f^{*}$ à \unsure{isom|iso} unique près. Mais en général, il ne suffit pas de se donner un tel $f_{*}$, pour pouvoir affirmer qu'il correspond à un $f$, i.e qu'il admette un adjoint à droite ; pour ceci, il faudrait qu'on sache que les foncteurs $M\unsure{\text{---}|\longrightarrow}\mathrm{Ens}$ qui commutent aux petites $\underset{\longleftarrow}{\lim}$ sont représentables, ce qui sera le cas si $M$ admet une petite sous-catégorie \uline{co}génératrice, i.e. si $M^{\mathrm{o}}$ est également une catégorie pseudo-topos.

Il est sans doute judicieux de procéder comme dans le cas des topos, et d'appeler \uline{morphisme de pseudo-topos} $f:\unsure{X|\mathcal X}\to\unsure{Y|\mathcal Y}$ un \unsure{quadruple|quadruplet}

\[
(f^{*},f_{*},\alpha,\beta)
\tag*{(4.1.3)}
\]
% archive sheet number: 135
\scanpage{137}
\mspage{53}

de deux foncteurs adjoints $(f^{*},f_{*})$ et des morphismes d'adjonction. Quand $f^{*}$ est une équivalence, ce qui équivaut à $f_{*}$ une équivalence, ou $\alpha$ et $\beta$ des iso, on dira qu'on a une \uline{équivalence de pseudo-topos}, et à toutes fins pratiques, on pourra alors regarder $\unsure{X|\mathcal X}$ et $\unsure{Y|\mathcal Y}$ comme des objets essentiellement identiques.

\deleted{\illegible{1}} Pour deux pseudo-topos $X,Y$, j'ai envie d'introduire (comme dans le cas des topos) un pseudo-topos produit

\[
X\times Y,
\tag*{(4.1.4)}
\]

% slanted
\margin{Il apparaît finalement que ce qu'on construit n'est \uline{pas} un 2-produit \unsure{! (dans la|(dans la} 2-catégorie envisagée.) \deleted{\unsure{donc}} il vaut mieux le noter par un signe tel que $X\boxtimes Y$. Mais si $X,Y$ sont deux topos, $X\boxtimes Y$ est un 2-produit dans la 2-catégorie des topos.}

ayant les propriétés d'un 2-produit dans la 2-catégorie des pseudo-topos. \unsure{mais|Mais} précisons que par définition

\[
\begin{aligned}
\underline{\operatorname{Hom}}_{\mathrm{ps\,top}}(X,Y)
&\simeq\underline{\operatorname{Hom}}_{!}(F(Y),F(X))^{\mathrm{o}},\\
&\simeq\underline{\operatorname{Hom}}^{!'}(F(X),F(Y)),
\end{aligned}
\tag*{(4.1.5)}
\]

où l'\deleted{exposant}\inserted{indice} ! \inserted{(resp. l'exposant !)} indique la sous-catégorie pleine de la catégorie $\underline{\operatorname{Hom}}$ envisagée, formée des foncteurs qui commutent aux \inserted{petites} $\underset{\longrightarrow}{\lim}$ (resp. aux petites
% archive sheet number: 136
\scanpage{138}
\mspage{54}

$\underset{\longleftarrow}{\lim}$), et \inserted{l'accent dans} le $\underline{\operatorname{Hom}}^{!'}$ dans le dernier \unsure{mb|membre} indique qu'on exige de plus que le foncteur en question ait un adjoint \uline{à gauche}.

Soient donc \deleted{$M,N$}

\[
M=F(\unsure{X|\mathcal X}),\quad N=F(\unsure{Y|\mathcal Y})
\tag*{(4.1.6)}
\]

et soit $\unsure{Z|\mathcal Z}$ un troisième ps. topos, décrit par

\[
P=F(Z),
\tag*{(4.1.7)}
\]

On veut expliciter, à équivalence près, la catégorie produit $\underline{\operatorname{Hom}}_{\mathrm{ps\,top}}(Z,X)\times\underline{\operatorname{Hom}}_{\mathrm{ps\,top}}(Z,Y)$, i.e.

\[
\underline{\operatorname{Hom}}_{!}^{\deleted{\unsure{!}}}(M,P)^{\mathrm{o}}
\times\operatorname{Hom}_{!}(N,Q)^{\mathrm{o}}
\tag*{(4.1.8)}
\]

\deleted{Si $(f^{*},g^{*})$ est un objet de cette catégorie, je \unsure{dois|désire} leur associer un bifoncteur}

\deleted{\[
\deleted{\unsure{f\times g}}\,f^{*}\boxtimes g^{*}:M\times N\longrightarrow P
\tag*{(4.1.9)}
\]}

\deleted{par}

\deleted{\[
(f^{*}\boxtimes g^{*})(x,y)=
\underbrace{\unsure{g^{*}|f^{*}}(\unsure{a|x})\times f^{\unsure{*|\alpha}}(y)}_{\substack{\text{produit cartésien}\\\text{dans }P}}
\tag*{(4.1.10)}
\]}

Je vais introduire pour cela une catégorie
% archive sheet number: 137
\scanpage{139}
\mspage{55}

\[
M\boxtimes N
\tag*{(4.1.9)}
\]

par la formule

\deleted{\[
\begin{aligned}
M\boxtimes N &\unsure{\simeq|=} \text{sous-catégorie pleine de}\\
&\underline{\operatorname{Hom}}^{!}(M^{\mathrm{o}},N)
\end{aligned}
\tag*{(\unsure{4.1.10})}
\]}

\[
M\boxtimes N\simeq\underline{\operatorname{Hom}}^{!!}(M^{\mathrm{o}}\times N^{\mathrm{o}},\mathrm{Ens})
\tag*{(4.1.10)}
\]

où le signe $\underline{\operatorname{Hom}}^{!!}$ désigne la sous-catégorie pleine \unsure{dans|de} la catégorie $\underline{\operatorname{Hom}}$, formée des bi-contrafoncteurs \deleted{\unsure{$F$|$B$}}

\[
F:M^{\mathrm{o}}\times N^{\mathrm{o}}\unsure{\longrightarrow|\text{---}}\mathrm{Ens}
\tag*{(4.1.10.1)}
\]

tels que pour tout $x$ dans $M$, le foncteur

\[
y\longmapsto F(\unsure{x|a},y):N^{\mathrm{o}}\unsure{\text{---}|\longrightarrow}\mathrm{Ens}
\tag*{(4.1.10.2)}
\]

\deleted{\unsure{transforme}} soit dans $\underline{\operatorname{Hom}}^{!}(N^{\mathrm{o}},\mathrm{Ens})$. i.e transforme \inserted{petites} $\underset{\longrightarrow}{\lim}$ en $\underset{\longleftarrow}{\lim}$, et que, symétriquement, pour \unsure{tt|tout} $y$ dans $N$, le foncteur

\[
x\longmapsto F(x,y):M^{\mathrm{o}}\text{---}\mathrm{Ens}
\tag*{(4.1.10.3)}
\]

soit dans $\underline{\operatorname{Hom}}^{!}(M^{\unsure{\mathrm{o}|b}},\mathrm{Ens})$. Mais cela signifie aussi que ces foncteurs sont \uline{représentables} ; donc on trouve aussi

\[
M\boxtimes N\xleftarrow{\approx}\underline{\operatorname{Hom}}^{!}(N^{\mathrm{o}},M)
\tag*{(4.1.11)}
\]

et

\[
M\boxtimes N\xleftarrow{\approx}\underline{\operatorname{Hom}}^{!}(M^{\mathrm{o}},N)
\tag*{(4.1.12)}
\]
% archive sheet number: 138
\scanpage{140}
\mspage{56}

En utilisant (\unsure{dans} chacune de ces formules) des deux côtés les équivalences de catégories

\[
\left\{\begin{aligned}
M&\xrightarrow{\approx}\underline{\operatorname{Hom}}^{!}(M^{\unsure{\mathrm{o}|b}},\mathrm{Ens})\\
N&\xrightarrow{\approx}\underline{\operatorname{Hom}}^{!}(N^{\mathrm{o}},\mathrm{Ens})
\end{aligned}\right.
\tag*{(4.1.13)}
\]

On peut dire \unsure{aussi|qu'ici} que la catégorie $M\overline{\boxtimes}N$ est équivalente à celle des \uline{couples} de foncteurs

\[
\left\{\begin{aligned}
\varphi &: M^{\mathrm{o}}\longrightarrow N\\
\psi &: N^{\mathrm{o}}\unsure{\longrightarrow|\text{---}}M
\end{aligned}\right.
\tag*{(4.1.\unsure{14|13})}
\]

\deleted{avec $\varphi\simeq\unsure{\psi}$} \deleted{\unsure{avec deux}} \deleted{\unsure{et}} munis \unsure{de|des} deux morphismes d'adjonction

\[
\left\{\begin{aligned}
\psi\varphi^{\mathrm{o}}&\xleftarrow{\alpha}\mathrm{id}_{M}\\
\varphi\psi^{\mathrm{o}}&\xleftarrow{\beta}\mathrm{id}_{N}
\end{aligned}\right.
\tag*{(4.1.15)}
\]

satisfaisant les deux compatibilités habituelles, exprimant l'adjonction de $\varphi^{\mathrm{o}}$ et de $\psi$ (ou, de façon équivalente, de $\psi^{\mathrm{o}}$ et de $\varphi$). Le foncteur

\[
(\varphi,\psi)\longmapsto\varphi\qquad(\text{resp. }(\varphi,\psi)\longmapsto\psi)
\]

donne une équivalence de la catégorie de ces couples avec $\deleted{\unsure{M}}\underline{\operatorname{Hom}}^{!}(M^{\mathrm{o}},N)$ (resp avec $\underline{\operatorname{Hom}}^{!}(N^{\mathrm{o}},M)$), de sorte qu'on a un carré commutatif d'équivalences \unsure{canoniques} de foncteurs :
% archive sheet number: 139
\scanpage{141}
\mspage{57}

\diagram{\begin{tikzcd}[ampersand replacement=\&,column sep=small,row sep=small]
\& \& M\overline{\boxtimes}N
   \arrow[dl,"\approx"'] \arrow[dr,"\approx"] \& \\
\text{(4.1.16)} \& \underline{\operatorname{Hom}}^{!}(M^{\mathrm{o}},N)
   \arrow[dr,"\approx"] \& \&
   \underline{\operatorname{Hom}}(N^{\unsure{\mathrm{o}|b}},M)\arrow[dl,"\approx"']\\
\& \& \begin{gathered}
 M\boxtimes N\\[-2pt]
 \overset{\mathrm{def}}{=}\underline{\operatorname{Hom}}^{!!}
 (M^{\mathrm{o}}\times N^{\mathrm{o}},\mathrm{Ens}).
 \end{gathered} \&
\end{tikzcd}}

\deleted{Je vais définir des foncteurs}

Je dis que $M\boxtimes N$ \uline{est une catégorie pseudo-topos}. Il est clair, certes, qu'elle admet des petites $\underset{\longleftarrow}{\lim}$, qui se calculent argument par argument dans $\underline{\operatorname{Hom}}(M^{\mathrm{o}}\times N^{\mathrm{o}},\mathrm{Ens})$, ou au choix dans $\underline{\operatorname{Hom}}^{!}(M^{\mathrm{o}},N)$, ou $\underline{\operatorname{Hom}}^{!}(N^{\mathrm{o}},M)$. Comment construire des petites $\underset{\longrightarrow}{\lim}$ ? Il n'est pas question d'invoquer un argument de cogénérateurs. Il faut d'ailleurs trouver aussi une petite sous-catégorie strictement génératrice. Visiblement, il y a \unsure{qque|qqe} chose de non trivial à démontrer !

4.2. \inserted{\uline{Cafouillages} sur $M\boxtimes N$.}

Revenons à la situation générale d'un pseudo-topos $\unsure{X|\mathcal X}$, décrit par la catégorie pseudo-topos $M$. Si $N$ est une catégorie \deleted{\unsure{satisfaisant|stable} \unsure{quelques|par} petites $\unsure{\underset{\longleftarrow}{\lim}|\underline{\lim}}$} définit la catégorie
% archive sheet number: 140
\scanpage{142}
\mspage{58}

\[
\operatorname{Fais}(\unsure{X|\mathcal X},N)\quad\unsure{\text{--}|\text{ou}}\quad F(X,N)
\tag*{\deleted{(\unsure{4.2.1})}}
\]

des faisceaux \unsure{sur|en} $X$ à valeurs dans $N$) par

\[
F(X,N)\overset{\mathrm{def}}{=}\underline{\operatorname{Hom}}^{!}(M^{\mathrm{o}},N).
\tag*{(4.2.1)}
\]

On a donc un foncteur pleinement fidèle

\[
F(X,N)\overset{\mathrm{def}}{=}\underline{\operatorname{Hom}}^{!}(M^{\mathrm{o}},N)
\hookrightarrow\underline{\operatorname{Hom}}^{!}(N^{\mathrm{o}},\overline M),
\tag*{(4.2.2)}
\]

où $\overline M$ désigne la sous-catégorie de $M^{\wedge}$ (\uline{NB} $M^{\wedge}$ n'est pas en général une $\unsure{U|\mathcal U}$-catégorie, \unsure{qu'à cela ne tienne|qui à cela ne tienne}) formée des foncteurs \uline{représentables} \deleted{\unsure{les}} $M^{\mathrm{o}}\to\mathrm{Ens}$ (laquelle est une $U$-catégorie, équivalente à $M$ par l'\unsure{inclusion|induction} canonique $M\unsure{\hookrightarrow|\text{---}}\overline M$). Donc on a à isom unique près une factorisation de (4.2.2) en

\[
F(X,N)\overset{\text{pl. fid.}}{\unsure{\hookrightarrow|\longleftrightarrow}}
\underline{\operatorname{Hom}}^{!}(N^{\mathrm{o}},M),
\tag*{(4.2.3)}
\]

l'image essentielle étant formée des foncteurs

\[
\psi:N^{\mathrm{o}}\unsure{\text{---}|\longrightarrow}M
\]

tels que pour \unsure{tout|tt} $x$ dans $M$,

\[
y\longmapsto\operatorname{Hom}(y,\psi(x)):N^{\mathrm{o}}\text{---}\mathrm{Ens},
\]

soit représentable, i.e. des foncteurs qui
% archive sheet number: 141
\scanpage{143}
\mspage{59}

admettent un \uline{adjoint à gauche}. Quand $N$ lui-$=$ est une catégorie pseudo-topos, cette condition est automatique (elle résulte de l'hypothèse que $\psi$ commute aux petites $\underset{\longleftarrow}{\lim}$). Donc dans ce cas, (4.2.3) est une équivalence de catégories.

Soit maintenant $S$ \inserted{($S$ comme ``site'')} une \uline{petite} sous-catégorie de $M$, génératrice par épimorphismes stricts. On sait que cela signifie que le foncteur composé

\[
M\xrightarrow{\approx}
\deleted{\unsure{M^{\wedge}}}
\deleted{\underline{\operatorname{Hom}}^{\unsure{? !|!}}(M^{\mathrm{o}},\mathrm{Ens})}
\inserted{\underline{\operatorname{Hom}}^{!}(M^{\mathrm{o}},\mathrm{Ens})}
\underset{\text{restr. à }S}{\longrightarrow}
S^{\wedge}\;(\overset{\mathrm{def}}{=}\underline{\operatorname{Hom}}(S^{\mathrm{o}},\mathrm{Ens}))
\tag*{(4.2.4)}
\]

est pleinement fidèle, et ceci implique, plus généralement, que pour \unsure{tte|toute} \deleted{\unsure{une}} catégorie $N$, le foncteur restriction

\[
F(\unsure{X|\mathcal X},N)=\underline{\operatorname{Hom}}^{!}(M^{\mathrm{o}},N)
\overset{\text{pl. fid.}}{\hookrightarrow}\underline{\operatorname{Hom}}(S^{\mathrm{o}},N)
\tag*{(4.2.5)}
\]

est pleinement fidèle. En effet, le deuxième membre, via le plongement pl. fidèle

\[
N\overset{\text{pl. fid}}{\hookrightarrow}\underline{\operatorname{Hom}}^{!}(N^{\mathrm{o}},\mathrm{Ens})\hookrightarrow N^{\wedge},
\]

\deleted{s'identifie} se plonge dans

\[
\underline{\operatorname{Hom}}(S^{\mathrm{o}},\underline{\operatorname{Hom}}^{!}(N^{\mathrm{o}},\mathrm{Ens}))
\simeq \operatorname{Hom}^{!}(N^{\mathrm{o}},S^{\wedge})
\]

et la pleine fidélité de (4.2.5) équivaut alors à celle du foncteur composé

\[
\begin{aligned}
F(X,N)\overset{\mathrm{def}}{=}\underline{\operatorname{Hom}}^{!}(M^{\mathrm{o}},N)&\longrightarrow
\operatorname{Hom}^{!}(N^{\mathrm{o}},S^{\wedge})\\
&\simeq\underline{\operatorname{Hom}}(S^{\mathrm{o}},\underline{\operatorname{Hom}}^{!}(N,\mathrm{Ens})),
\end{aligned}
\tag*{(4.2.6)}
\]
% archive sheet number: 142
\scanpage{144}
\mspage{60}

Or on a commutativité dans le carré

\diagram{\begin{tikzcd}[ampersand replacement=\&,column sep=large,row sep=large]
\& F(\unsure{X|\mathcal X},N)\overset{\mathrm{def}}{=}\underline{\operatorname{Hom}}^{!}(M^{\mathrm{o}},N)
 \arrow[r,"{(4.2.5)}"] \arrow[d,hook,"{(4.2.3)}"',"{\text{pl. fid}}"]
 \& \underline{\operatorname{Hom}}(S^{\mathrm{o}},N)\arrow[d,hook,"{\text{pl. fid}}"]\\
\text{(4.2.7)} \& \underline{\operatorname{Hom}}^{!}(N^{\mathrm{o}},M)
 \arrow[r,hook,"{(4.2.4)}","{\text{pl. fid}}"']
 \& \begin{gathered}\underline{\operatorname{Hom}}^{!}(N^{\mathrm{o}},S^{\wedge})\\
 \simeq\underline{\operatorname{Hom}}(S^{\mathrm{o}},\underline{\operatorname{Hom}}^{!}(N^{\mathrm{o}},\mathrm{Ens}))\end{gathered}
\end{tikzcd}}

(où les flèches verticales sont des équivalences, si $N$ est lui-$=$ une cat. pseudo-topos).

\deleted{Ainsi : \uline{Proposition} 4.2.1 Soient $M\inserted{=F(X)},N\inserted{=F(Y)}$ deux catégories pseudo-topos, \inserted{et $S$ une \deleted{\unsure{sous}} \unsure{catégorie|sous-catégorie} str. gén. de $M$.} Alors $F(X,N)\overset{\mathrm{def}}{=}\underline{\operatorname{Hom}}^{!}(M^{\mathrm{o}},N)$}

\deleted{\deleted{\unsure{et}} $\simeq\underline{\operatorname{Hom}}^{!}(N^{\mathrm{o}},M)\approx M\boxtimes N$ se plonge \deleted{dans} par \deleted{le}\inserted{un} foncteur \inserted{-restriction} \uline{pleinement fidèle} en $\underline{\operatorname{Hom}}(S^{\mathrm{o}},N)\deleted{\unsure{\simeq}}\simeq\deleted{\unsure{F}}S^{\wedge}\boxtimes N$}

Notons d'autre part \deleted{\unsure{que}} que l'on a une équivalence

\[
\underline{\operatorname{Hom}}(S^{\mathrm{o}},N)\xleftarrow{\approx}
\underline{\operatorname{Hom}}^{!}(S^{\wedge\,\mathrm{o}},N)
\overset{\mathrm{def}}{=}F(\operatorname{Top}(S),N)
\tag*{(4.2.8)}
\]

où $\operatorname{Top}(S)$ est le pseudo-topos \inserted{(en fait, un topos)} défini par la catégorie pseudo-topos $S^{\wedge}$. Pour \deleted{\unsure{catégorie}} \inserted{ceci} notons qu'on a

\[
\underline{\operatorname{Hom}}^{!}(S^{\wedge\,\mathrm{o}},N)\xrightarrow{\sim}
\underline{\operatorname{Hom}}_{!}(S^{\wedge},N^{\mathrm{o}})^{\mathrm{o}}
\]

et d'autre part le foncteur restriction à $S\subset S^{\wedge}$ donne

\[
\underline{\operatorname{Hom}}_{!}(S^{\wedge},N^{\mathrm{o}})\xrightarrow{\approx}
\underline{\operatorname{Hom}}(S,N^{\mathrm{o}})
\]

d'où
% archive sheet number: 143
\scanpage{145}
\mspage{61}

\[
\underline{\operatorname{Hom}}(S^{\wedge\,\mathrm{o}},N)\xrightarrow{\approx}
(\underline{\operatorname{Hom}}(S,N^{\mathrm{o}}))^{\mathrm{o}}\xrightarrow{\sim}
\underline{\operatorname{Hom}}(S^{\mathrm{o}},N),
\]

ce qui n'est autre que (4.2.8). Ainsi

\uline{Proposition} 4.2.1. Soient $M=F(\unsure{X|\mathcal X})$ \inserted{\unsure{munie|muni} d'une sous-cat. pleine str. génératrice $S$,} une catégorie pseudo-topos, $N$ une $\unsure{\mathcal U|U}$-catégorie quelconque, alors le foncteur ``restriction à $S$'' suivant est pleinement fidèle

\[
F(X,N)\overset{\mathrm{def}}{=}\underline{\operatorname{Hom}}^{!}(M^{\mathrm{o}},N)
\longrightarrow\underline{\operatorname{Hom}}(S^{\mathrm{o}},N)
\xleftarrow{\approx}\begin{gathered}\underline{\operatorname{Hom}}^{!}(S^{\wedge\,\mathrm{o}},N)\\
\mathrel{\Vert}\,\mathrm{def}\\F(\operatorname{Top}(S),N),\end{gathered}
\tag*{(4.2.1.\deleted{\unsure{2|3}}1)}
\]

et il s'insère dans le carré commutatif (4.2.7) de foncteurs pleinement fidèles.

\uline{Corollaire} 4.2.2. Supposons que $N\inserted{=F(Y)}$ soit également une catégorie pseudo-topos, munie d'une sous-catégorie pleine $T$ str. génératrice. Alors

\[
F(X,N)\simeq F(Y,M)\simeq
\underbrace{\underline{\operatorname{Hom}}^{!!}(M^{\mathrm{o}}\times N^{\mathrm{o}},\mathrm{Ens})}_{=M\boxtimes N},
\]

et le foncteur restriction

\[
M\boxtimes N\overset{\mathrm{def}}{=}
\underline{\operatorname{Hom}}^{!!}(M^{\mathrm{o}}\times N^{\mathrm{o}},\mathrm{Ens})\longrightarrow
\begin{gathered}\underline{\operatorname{Hom}}(S^{\mathrm{o}}\times T^{\mathrm{o}},\mathrm{Ens})\\
\rotatebox{90}{\(\simeq\)}\\(S\times T)^{\wedge}\\\mathrel{\Vert}\,\mathrm{def}\\F(\operatorname{Top}(S\times T))\end{gathered}
\tag*{(4.2.2-1)}
\]

est pleinement fidèle.

La première assertion a été vue en 4.1. Pour la pleine fidélité de (4.2.2.1), on note que ce \deleted{\unsure{En effet, c'est|La flèche, c'est}} \unsure{dernier|foncteur} s'interprète, à équiv. de cat. près, comme le composé du foncteur pleinement fidèle (4.2.1.1), et du foncteur
% archive sheet number: 144
\scanpage{146}
\mspage{62}

similaire relatif à $N$ (\inserted{donc} lui aussi pl. fid.)

\[
\deleted{\unsure{F(X|F(N}}F(\unsure{Y|\mathcal Y},S^{\wedge})\overset{\mathrm{def}}{=}
\underline{\operatorname{Hom}}^{!}(N^{\wedge},S^{\wedge})
\xrightarrow{\text{restr. à }T}\begin{gathered}\underline{\operatorname{Hom}}(T^{\mathrm{o}},S^{\wedge})\\
\rotatebox{90}{\(\simeq\)}\\\underline{\operatorname{Hom}}(S^{\mathrm{o}}\times T^{\mathrm{o}},\mathrm{Ens}),\end{gathered}
\tag*{(4.2.2.2)}
\]

compte tenu de l'\unsure{équivalence|identification} de symétrie

\[
F(Y,S^{\wedge})\overset{\mathrm{def}}{=}
\underline{\operatorname{Hom}}^{!}(N^{\deleted{\unsure{\wedge}}\mathrm{o}}, S^{\wedge})\simeq
\underline{\operatorname{Hom}}^{!}(S^{\wedge\,\mathrm{o}},N).
\tag*{(4.2.2.3)}
\]

Au total, quand $M,N$ sont tous deux des cat. pseudo-topos, munies des \inserted{ss-}cat. strictement génératrices $S$ resp. $T$, on a un plongement pleinement fidèle canonique

\[
M\boxtimes N\overset{\unsure{\rho|\varrho}}{\lhook\joinrel\longrightarrow}(S\times T)^{\wedge},
\tag*{(4.2.9)}
\]

qu'on interprétera par la suite comme $(f^{M}_{S}\times g^{N}_{T})_{*}$, où \deleted{$\unsure{f^{M}_{S}:X|f_{S}:X}$}

\[
\left\{\begin{aligned}
f^{M}_{S}&: \unsure{X|\mathcal X}\longrightarrow\operatorname{Top} S\\
f^{N}_{\deleted{\unsure{S}}T}&: Y\longrightarrow\operatorname{Top}\deleted{\unsure{S}}T
\end{aligned}\right.
\tag*{(4.2.10)}
\]

sont les \unsure{homorphismes|homomorphismes} de pseudo-topos \deleted{\unsure{dont ils les}}, par les plongements pl. fidèles (\inserted{foncteurs} images directes)

\[
\left\{\begin{aligned}
(f^{M}_{S})_{*}&:M\lhook\joinrel\longrightarrow S^{\wedge}\\
(f^{N}_{T})_{*}&:N\lhook\joinrel\longrightarrow T^{\wedge},
\end{aligned}\right.
\tag*{(4.2.11a)}
\]

lesquels admettent bel et bien des adjoints à gauche
% archive sheet number: 145
\scanpage{147}
\mspage{63}

\[
\left\{\begin{aligned}
(f^{M}_{S})^{*}&: S^{\wedge}\longrightarrow M\\
(f^{N}_{T})^{\unsure{\prime|*}}&: T^{\wedge}\longrightarrow N,
\end{aligned}\right.
\tag*{(4.2.12)}
\]

à savoir les prolongements canoniques (\unsure{qui commutent|par commutation} aux petites $\unsure{\underrightarrow{\lim}|\underline{\lim}}$) des foncteurs d'inclusion $S\lhook\joinrel\longrightarrow M$, $T\lhook\joinrel\longrightarrow N$.

La question essentielle, à présent, c'est évidemment de construire un \deleted{\unsure{$f$}} adjoint à gauche pour le foncteur (4.2.9). Il revient au $=$, pour $(s,t)$ dans $S\times T$, de définir un objet

\[
s\boxtimes t\in\operatorname{Ob} M\boxtimes N
\tag*{(4.\deleted{\unsure{1}}2.13)}
\]

tel que l'on ait, pour $F$ variable dans $M\boxtimes N$, un isom. fonctoriel en $F$

\[
\underbrace{\unsure{\rho|\varrho}(F)(s,t)}_{\substack{\Vert\,\mathrm{def}\\ \deleted{\unsure{\mathrm{def}}}F(s,t)}}\simeq\operatorname{Hom}(s\boxtimes t,F):
\]

Donc il faut prouver ceci

% slanted
\margin{Pas prouvé, peut-être faux !}

\uline{Lemme} 4.2.3. Soient $M,N$ deux catégories pseudo-topos, et $\deleted{\unsure{x}}s\in\operatorname{Ob}M$, $\deleted{\unsure{y}}t\in\operatorname{Ob}N$, et considérons sur $M\boxtimes N\overset{\mathrm{def}}{=}\underline{\operatorname{Hom}}^{!!}(M^{\mathrm{o}}\times N^{\mathrm{o}},\mathrm{Ens})$ le foncteur $F\longmapsto F(\deleted{\unsure{x}}s,\deleted{\unsure{y}}t)$. Ce foncteur est \uline{représentable} (par un objet que je note $\deleted{\unsure{x}}s\boxtimes\deleted{\unsure{y}}t$).

L'idée naturelle serait de définir $\deleted{\unsure{x}}s\boxtimes\deleted{\unsure{y}}t$ par

\[
(\deleted{\unsure{x}}s\boxtimes\deleted{\unsure{y}}t)(\deleted{\unsure{s}}x,y)
=\deleted{\unsure{x}}s(\deleted{\unsure{s}}x)\times\deleted{\unsure{y}}t(y)
\unsure{\equiv|=}\operatorname{Hom}_{M}(x,s)\times\operatorname{Hom}_{N}(y,t)
\]
% archive sheet number: 146
\scanpage{148}
\mspage{64}

ce qui définit bien un foncteur

\[
M^{\mathrm{o}}\times N^{\mathrm{o}}\longrightarrow\mathrm{Ens}.
\]

Malheureusement, pour un des arguments \inserted{fixé}, disons pour $y$ fixé, \deleted{\unsure{c'est}} ce foncteur \deleted{\unsure{un foncteur}}\inserted{n'est} \uline{pas} \deleted{\unsure{limite}} compatible aux limites projectives, p.ex. il ne transforme pas l'objet final de $N^{\mathrm{o}}$ (correspondant à l'objet initial de $N$) en objet final de $\mathrm{Ens}$, mais en \deleted{$\unsure{\operatorname{Hom}(}$} $t(y)=\operatorname{Hom}_{N}(y,t)$. Néanmoins, je pense que

% slanted
\margin{également faux, cf. plus bas}

$s\boxtimes t$ doit bel et bien se définir comme produit de deux objets de $M\boxtimes N$, dépendant fonctoriellement de \unsure{$s$ resp. de $t$|$s$ resp de $t$}. On est donc \unsure{amené|ramené} dès à présent, qu'on le veuille ou non, \deleted{\unsure{à ce que}}, \unsure{à définir|à la question|à la définition} des foncteurs canoniques, commutant aux petites $\unsure{\underrightarrow{\lim}|\underline{\lim}}$

\[
M\longrightarrow M\boxtimes N\longleftarrow N,
\tag*{(4.2.\unsure{14|13})}
\]

% slanted
\margin{\deleted{\unsure{Le}} Cas où $M,N$ sont des catégories-topos nous montre que ces foncteurs doivent être décrits \unsure{comme} $x\longmapsto x\boxtimes e_{N}$, et $y\longmapsto e_{M}\boxtimes y$. Mais il n'est pas clair \unsure{qu'on ait|qu'il y ait}, en dehors du cas des topos, et en admettant l'existence des $x\boxtimes y$, la formule (valable pour les topos) $x\boxtimes y\simeq(x\boxtimes e_{N})\times(e_{M}\boxtimes y)$}

(lesquels seront les foncteurs images inverses des projections canoniques

\diagram{\begin{tikzcd}[ampersand replacement=\&,column sep=large,row sep=small]
\&\unsure{X|\mathcal X}\\
X\times Y\arrow[ur]\arrow[dr]\&\\
\& Y
\end{tikzcd}}

).

4.3. \uline{L'opération} $M\boxtimes N$ \uline{et l'axiome d'accessibilité}.

J'ai passé la journée hier, essayant d'établir l'existence des $x\boxtimes y$ dans $M\boxtimes N$, sans y parvenir -- du moins pas sans faire \deleted{\unsure{des}} une hypothèse supplémentaire d'accessibilité,
% archive sheet number: 147
